% !PRINTED 432
\weberpage{432}
%<<CONT
and from
\weq{10}{p\,\alpha'\,\delta-\beta\,\gamma = 1}
there then follows
\weq{11}{y^2+m\,x^2 = 4.}

The value $x=0$ is to be excluded, because otherwise $\beta=0$ would have to hold, which contradicts equation (10), and we may, without loss of generality, assume $x$ to be positive. Since moreover $m$ must be either $\equiv 0$ or $\equiv 3$ with respect to the modulus 4, there remain for the satisfaction of (11) only the following two possibilities:
\[
1.\quad m=3,\quad x=1,\quad y=\pm 1.
\]

Since for the discriminant $-3$ there is only one class of forms, we may assume $A=1,\,B=1,\,C=1$, that is, we may set for $\omega$ the imaginary third root of unity $e^{\frac{2\pi i}{3}}$, and we then obtain among the quantities (3) three equivalent ones by taking $y=+1$ and $=-1$: % [sic] The German print omits the variable y before the second equals sign.
\weq{12}{p\,\omega,\quad \frac{\omega}{p},\quad \frac{\omega+1}{p},}
as is also recognised from the equations
\[
\frac{\omega}{p} = \frac{-1}{p+p\,\omega}, \quad \frac{\omega+1}{p} = \frac{-1}{p\,\omega}
\]
from which the equivalence of the three quantities (12) is evident.
\[
2.\quad m=4,\quad x=1,\quad y=0.
\]

Because here too there is only one class of forms, we may assume $B=0,\,A=C=1$, that is, $\omega=i$, and we find the two equivalent values:
\[
p\,\omega,\quad \frac{\omega}{p},
\]
as also follows from the equation
\[
p\,\omega = \frac{-p}{\omega}
\]

Further, let two of the values (3)
\[
\frac{\omega+c}{p}, \quad \frac{\omega+c'}{p}
\]
be equivalent, so that an equation holds:
\weq{13}{\frac{\omega+c}{p} = \frac{\gamma\,p+\delta\,(\omega+c')}{\alpha\,p+\beta\,(\omega+c')}, \quad \alpha\,\delta-\beta\,\gamma = 1,}
%CONT>>
